\input{./._relpath-to-latexroot.ltx}
\documentclass[\latexroot/\projectname]{subfiles}
\input{\latexroot/@resources/subfile-setup.ltx}

\begin{document}

\section{Experiments}
\whenintegrated{\label{sec:experiments}}

To study the implications of a reversal in capital flows on the model economy, we run two separate experiments. For the first, we consider an experiment inspired by the Hungarian current account expansion and reversal of the early and late 2000s. Starting from an initial steady state with low external liabilities, the economy experiences a sustained expansion in the current account that then stabilizes, followed by an unexpected reversal.

In the baseline experiment we assume that the model economy follows a floating exchange rate regime with all debt foreign-currency denominated, consistent with the Hungarian experience.

\begin{figure}[htbp]
  \centering
  \includegraphics[width=0.8\textwidth]{Figures/deFerra2020_fig1_credit_supply.png}
  \caption{Foreign credit supply shocks.}
  \whenintegrated{\label{fig:credit_supply}}
\end{figure}

\section{Sudden stop}
\whenintegrated{\label{sec:suddenstop}}

The main experiment of interest is studying the reversal of foreign credit supply under the baseline scenario of a flexible exchange rate and foreign-currency denominated debt---the empirically relevant case for Hungary during the Global Financial Crisis, where 69\% of household debt was denominated in foreign currency.

Under a flexible exchange rate policy, consumption of home and foreign goods by domestic consumers contracts sharply because of the unexpected withdrawal of foreign credit and the fall in wealth due to the debt revaluation. We find, using a utilitarian metric, that under the flexible-exchange rate regime, welfare falls by 0.81\% in consumption equivalent variation (CEV) terms.

More than 92\% of households would have benefited from the fixed-exchange rate policy. Despite the fact that the fixed exchange rate generates a domestic recession---due to the presence of nominal rigidities---the fixed exchange rate policy is preferred by the vast majority of households. Our framework therefore generates a rationale for ``fear of floating'' even in the absence of a contractionary devaluation.

% Bibliography inclusion handled automatically by document structure
\smartbib

\end{document}
